\clearpage
\begin{figure}[t!]
\huge
\begin{center}
    \textbf{Supplementary Material}
\end{center}
\end{figure}

\section{Additional Proofs}

\subsection{Proof of Proposition~\ref{prop:linear=cos}}\label{app:proof:prop:linear=cos}


\begin{proof}
Let $n=\med(m, k; \fcos)$ and there exists $X=\{x_1,\dots,x_m\}\subset\R^{n}\setminus\{0\}$ \footnote{If there is $x_i = 0$, the proof is basically the same, we omit this situation for convenience.}that is $k$-shattered by $\fcos$.
Define the radial projection $\rho:\R^{n}\setminus\{0\}\to S^{n-1}$ by $\rho(x)=\frac{x}{\|x\|_2}$, and let
$Y=\{\rho(x_i)\}_{i=1}^m\subset S^{n-1}$.

By $k$-shattering, for each $S\subseteq X$ with $|S|\le k$ there exist $\vw_S\in\R^{n}\setminus\{0\}$ such that
\[
x\in S \Rightarrow r_S-\frac{\langle \vw_S,x\rangle}{\|\vw_S\|\|x\|} \le 0,
\qquad
x\in X\setminus S \Rightarrow r_S-\frac{\langle \vw_S,x\rangle}{\|\vw_S\|\|x\|} > 0.
\]

For $z\in S^{n-1}$ define the affine functional
\[
f_S(z)\;:=\;\big\langle \tfrac{-\vw_S}{\|\vw_S\|},\, z \big\rangle + r_S
\ \in\ \fln.
\]

Then, for every $x\in X$ we have
\[
f_S(\rho(x)) \;=\; r_S-\frac{\langle \vw_S,x\rangle}{\|\vw_S\|\|x\|},
\]
so the inequalities above becomes
\[
y\in \rho(S)\Rightarrow f_S(y)\le 0,\qquad y\in Y\setminus \rho(S)\Rightarrow f_S(y)>0.
\]
Therefore we show that 
\[
\med(m, k; \fln) <= \med(m, k; \fcos)
\]

Conversely, let $n^\star=\mathrm{MED}(m,k;\fln)$. Then there exists $X=\{x_1,\dots,x_m\}\subset\R^{n^\star}$ such that for every $S\subseteq X$ with $|S|\le k$ there is $f_S(x)=\langle \vw_S,x\rangle+b_S$ satisfying
\[
x\in S \Rightarrow f_S(x)\le 0,\qquad x\in X\setminus S \Rightarrow f_S(x)>0.
\]

Consider the embedding $\phi:\R^{n^\star}\to S^{n^\star}\subset\R^{n^\star+1}$ given by
\[
\phi(x)\;=\;\frac{(x,1)}{\|(x,1)\|},
\]
where $(x,1)$ denotes the concatenation of $x$ and 1.

Let $Y=\{\phi(x_i)\}_{i=1}^m\subset S^{n^\star}$. For each $S$, set
\[
u_S\;=\;\frac{(\vw_S,b_S)}{\|(\vw_S,b_S)\|}\in S^{n^\star},\qquad t_S:=0,
\]
and define $g_S(y)=\tfrac{\langle u_S,y\rangle}{\|u_S\|\|y\|}-t_S=\langle u_S,y\rangle$ (since $\|u_S\|=\|y\|=1$ on $S^{n^\star}$).
Then for every $x\in\R^{n^\star}$,
\[
g_S(\phi(x))
=\left\langle \frac{(\vw_S,b_S)}{\|(\vw_S,b_S)\|},\,\frac{(x,1)}{\|(x,1)\|}\right\rangle
=\frac{\langle \vw_S,x\rangle+b_S}{\|(\vw_S,b_S)\|\cdot\|(x,1)\|}.
\]

Therefore, for $x\in X$ we have the exact shattering
\[
f_S(x)\le 0 \iff g_S(\phi(x))\le 0,\qquad
f_S(x)>0 \iff g_S(\phi(x))>0.
\]
\end{proof}

\subsection{Proof of Theorem~\ref{thm:medc-linear}, Theorem~\ref{thm:medc-cos}, and Theorem~\ref{thm:medc-l2}}\label{app:proof:thm:maed-linear}

Those three proofs are very similar, so they are grouped together.

\begin{proof}
Let $v_1, \dots, v_m \sim \mathcal{N}(0, I_n/n)$ in $\real^n$, our goals are to show that there is a positive probability such that 
for any subset $S$, $|S|\leq k$, $v_1\in S$ and $v_2\notin S$, the following inequalities holds:
\begin{align}
    \langle v_1, \sum_{u\in S} u \rangle & > \langle v_2 ,  \sum_{u\in S} u \rangle, & \text{for Theorem~\ref{thm:medc-linear}};\\
    \langle \frac{v_1}{\|v_1\|}, \sum_{u\in S} u \rangle & > \langle \frac{v_2}{\|v_2\|},  \sum_{u\in S} u \rangle, &\text{for Theorem~\ref{thm:medc-cos}};\\
    \|v_1 - \frac{1}{|S|} \sum_{u\in S} u \|_2^2 &< \|v_2 - \frac{1}{|S|} \sum_{u\in S} u \|_2^2, & \text{for Theorem~\ref{thm:medc-l2}}.
\end{align}

To show this, considering the concentration of the inner product of two independent random vectors $v_i$ and $v_j$ drawn from $\mathcal{N}(0, \frac{1}{n})$:
\begin{align}
\Pr\left(|\langle v_i, v_j\rangle| \geq \frac{1}{3k} \right) \leq 2\exp\left(-c\frac{n}{k^2}\right),
\end{align}
where $c$ is a constant.

Also, noticing that the norm of such vectors is concentrated as
\begin{align}
\Pr\left(|\|v_i\| - 1| \geq \frac{1}{3k} \right) \leq 2\exp\left(-c\frac{n}{k^2}\right),
\end{align}

\textbf{To prove theorem~\ref{thm:medc-linear}}, let $k=l$.
For $m$ vectors, take the union bound of (1) inner products of all possible pairs and (2) the norm of all vectors, then we derive the probability that any of the inner products (or the $\|v_i - 1\|$) is greater than $\frac{1}{3k}$, and then force it to be smaller than 1.
\begin{align}
	2 \left(m+\frac{m(m-1)}{2}\right) \exp\left(-c\frac{n}{k^2}\right) < 1.
\end{align}

The probability of the event where any inner products are smaller than $\frac{1}{3k}$, and norms are no diverge from 1 larger than $\frac{1}{3k}$ is positive when
\begin{align}
	n > C k^2 \log m.
\end{align}
where $C$ is a constant.

Under such conditions, 
\begin{align}
    &LHS = \langle v_1, \sum_{u\in S} u \rangle = \|v_1\| + \sum_{u\in S-\{v_1\}} \langle v_1, u\rangle > 1 - \frac{1}{3k} + (|S|-1) \frac{1}{3k} = 1 - \frac{|S|}{3k}, \\
    &RHS = \sum_{u\in S} \langle v_2, u\rangle < \frac{|S|}{3k}.
\end{align}
Because $|S| \leq k$, so $LHS > RHS$, Theorem~\ref{thm:medc-linear} is proved.

\textbf{To prove Theorem~\ref{thm:medc-cos}}
\begin{align}
    LHS &= \langle \frac{v_1}{\|v_1\|}, \sum_{u\in S} u \rangle > \frac{1 - |S| / 3k}{1 + 1/3k} = \frac{3k-|S|}{3k + 1} \\
    RHS &= \langle \frac{v_2}{\|v_2\|}, \sum_{u\in S} u \rangle < \frac{|S|}{3k(1-\frac{1}{3k})} = \frac{|S|}{3k-1}.
\end{align}
Then, 
\begin{align}
    LHS - RHS  > \frac{3k(3k - 2|S| -1)}{9k^2-1}.
\end{align}
Because $1 < |S| \leq k$, so $LHS - RHS > 0$ and Theorem~\ref{thm:medc-cos} is proved.

\textbf{To prove Theorem~\ref{thm:medc-l2}}
\begin{align}
    LHS & = \|v_1 - \frac{1}{|S|} \sum_{u\in S} u \|_2^2 = \|v_1\|_2^2 + \|\frac{1}{|S|} \sum_{u\in S} u \|_2^2 - \frac{2}{|S|}\langle v_1, \sum_{u\in S} u \rangle\\
    & < 1+\frac{1}{3k} - \frac{2}{|S|}(1-\frac{|S|}{3k})  + \|\frac{1}{|S|} \sum_{u\in S} u \|_2^2, \\
    RHS & = \|v_2 - \frac{1}{|S|} \sum_{u\in S} u \|_2^2 = \|v_2\|_2^2 + \|\frac{1}{|S|} \sum_{u\in S} u \|_2^2 - \frac{2}{|S|}\langle v_2, \sum_{u\in S} u \rangle \\
    & >  1-\frac{1}{3k} - \frac{2}{|S|}\frac{|S|}{3k} + \|\frac{1}{|S|} \sum_{u\in S} u \|_2^2.
\end{align}
Then,
\begin{align}
    LHS - RHS < \frac{2}{3k} + \frac{2}{3k} - \frac{2}{|S|}(1-\frac{|S|}{3k}) = \frac{2}{k} - \frac{2}{|S|}.
\end{align}
Because $|S| \leq k$, so $LHS - RHS < 0$ and Theorem~\ref{thm:medc-l2} is proved.
\end{proof}

\section{Simulation Code}\label{app:code}

\begin{lstlisting}
import torch
from torch import optim
from tqdm import trange
import itertools
import matplotlib.pyplot as plt
import numpy as np

class Trainer:
    def __init__(self, n, k):
        self.n = n
        self.k = k
        self.device="cuda" if torch.cuda.is_available() else "cpu"
        print("Using device", self.device)
        self.all_combinations = list(itertools.combinations(range(self.n), self.k))
        self.subset_indices_tensor = torch.tensor([list(s) for s in self.all_combinations], device=self.device)
        self.subset_excluded_tensor = torch.tensor(
            [
                [i for i in range(self.n) if i not in subset_indices]
                for subset_indices in self.all_combinations
            ], device=self.device
        )
        self.total_violations = len(self.all_combinations) * self.k * (self.n-self.k)

    def calculate_loss(self):
        # Convert batch_subset_indices to a tensor for efficient indexing
        # Calculate sums of vectors for each subset in the batch
        # Shape: (batch_size, d)
        subset_sums = self.vector_embeddings[self.subset_indices_tensor].mean(dim=1)
        # Calculate dot products of subset sums with all vectors
        # Shape: (batch_size, n)
        dot_products_all = torch.matmul(subset_sums, self.vector_embeddings.T)
        # Get dot products with vectors in the subset by indexing dot_products_all
        # Shape: (batch_size, k)
        dot_products_xi = torch.gather(dot_products_all, 1, self.subset_indices_tensor)
        # For each subset in the batch, calculate dot products with vectors NOT in the subset
        # We can create a mask to zero out elements within the subset
        # This is faster than iterating through the batch
        # Shape: (batch_size, n-k)
        dot_products_xj = torch.gather(dot_products_all, 1, self.subset_excluded_tensor)
        # Calculate differences and apply ReLU
        # We need to unsqueeze dot_products_xi to match dimensions for broadcasting
        # Shape: (batch_size, n-k) - (batch_size, k)
        differences = dot_products_xi.unsqueeze(1) - dot_products_xj.unsqueeze(2)
        total_loss = torch.relu(-differences).sum((-1, -2)).mean()
        num_violations = (differences < 0).sum().item()
        return total_loss, num_violations


    def train(self, d, num_epochs, learning_rate=0.1, patience=20):
        self.vector_embeddings = torch.randn(
            self.n, d, device=self.device, requires_grad=True
        )


        optimizer = optim.Adam(
            [self.vector_embeddings],
            lr=learning_rate,
        )
        scheduler = optim.lr_scheduler.OneCycleLR(
            optimizer=optimizer,
            max_lr=learning_rate,
            total_steps=num_epochs,
            pct_start=0.0,
        )

        min_violations = self.total_violations
        epochs_no_improve = 0

        with trange(num_epochs, desc=f"\t\t n={self.n}, k={self.k}, {d=}") as epoch_iterator:
            for epoch in epoch_iterator:
                optimizer.zero_grad()

                loss, violations = self.calculate_loss()
                loss.backward()
                optimizer.step()
                scheduler.step()


                epoch_loss = loss.item()
                epoch_iterator.set_postfix(
                    {
                        "n": self.n,
                        "k": self.k,
                        "loss": epoch_loss,
                        "#vio rate": violations / self.total_violations,
                        "min #vio": min_violations,
                        "epochs_no_improve": epochs_no_improve
                    }
                )

                if violations < min_violations:
                    min_violations = violations
                    epochs_no_improve = 0
                else:
                    epochs_no_improve += 1

                if violations == 0:
                    print("Early stopping: No violations found.")
                    break
                if epochs_no_improve >= patience:
                    print(f"Early stopping: No improvement in violations for {patience} epochs.")
                    break

        return min_violations


class Experiment:
    def __init__(self):
        self.search_paths = {}
        self.minimal_dimensions = {}

    def find_minimal_dimension(self, k, n_values, left0=0, num_epochs=100, learning_rate=0.1, patience=10):

        last_minimal = left0
        for n in n_values:
            print("#" * 10 + "new task" + "#" * 10)
            print(f"Finding minimal dimension for n={n}, k={k}")
            print("#" * 30)
            trainer = Trainer(n, k)

            self.search_paths[n] = []
            left, right = last_minimal+1, last_minimal + 40 # the range is set by prior to accelerate the convergence.
            minimal_d = n + 1

            while left <= right:
                mid = (left + right) // 2
                if mid == 0:
                    mid = 1

                print(f"\t>>>>Testing dimension d={mid}")

                violations = trainer.train(
                    d=mid,
                    num_epochs=num_epochs,
                    learning_rate=learning_rate / np.log2(n),
                    patience=patience
                )

                print(f"\t<<<<Violations for d={mid}: {violations}")

                self.search_paths[n].append({'dimension': mid, 'violations': violations})

                if violations == 0:
                    minimal_d = mid
                    right = mid - 1
                else:
                    left = mid + 1


            self.minimal_dimensions[n] = minimal_d if minimal_d <= n else -1
            last_minimal = minimal_d

            with open("minimal_dem_log.txt", "at") as f:
                f.write(f"minimal dimension @ {k=}&{n=} is {minimal_d}\n")
            print("#" * 10 + " Task Finished" + "#" * 10)
            print(f"minimal dimension @ {k=}&{n=} is {minimal_d}\n")
            print("#" * 30)

        return self.minimal_dimensions

    def plot_minimal_dimension_vs_n(self, k):
        if not self.minimal_dimensions:
            print("No experiment results to plot. Run find_minimal_dimension first.")
            return

        n_plot = list(self.minimal_dimensions.keys())
        d_plot = list(self.minimal_dimensions.values())

        plt.figure(figsize=(8, 6))
        plt.plot(n_plot, d_plot, marker='o', linestyle='-')
        plt.xlabel("Number of Vectors (n)")
        plt.ylabel("Minimal Dimension (d)")
        plt.xscale('log')
        plt.title(f"Minimal Dimension vs. Number of Vectors for k={k}")
        plt.grid(True)
        plt.show()

        plt.savefig("med_vs_n.png")

    def plot_violations_vs_d(self, k):
        if not self.search_paths:
            print("No experiment results to plot. Run find_minimal_dimension first.")
            return

        for n, path in self.search_paths.items():
            sorted_path = sorted(path, key=lambda x: x['dimension'])

            dimensions = [item['dimension'] for item in sorted_path]
            violations = [item['violations'] for item in sorted_path]

            plt.figure(figsize=(8, 6))
            plt.plot(dimensions, violations, marker='o', linestyle='-')
            plt.xlabel("Dimension (d)")
            plt.ylabel("Number of Violations")
            plt.title(f"Violations vs. Dimension for n={n}, k={k}")
            plt.grid(True)
            plt.show()


# 1. Define a list of n_values to experiment with
n_values = [5 * 2 ** (i) for i in [0, 1, 2, 3, 4, 5]]

# 2. Set a value for k
k = 2

# 3. Instantiate the Experiment class
experiment = Experiment()

# 4. Call the find_minimal_dimension method with a reasonable batch_size and limited epochs
#    Using a batch_size of 32 and num_epochs of 100 for initial testing.
print("Starting experiment with mini-batching and early stopping...")
minimal_dims = experiment.find_minimal_dimension(
    k, n_values, learning_rate=1, num_epochs=1000, patience=1000)

# 5. Print the minimal_dims dictionary
print("\n Minimal dimensions found:", minimal_dims)

# 6. Call the plot_minimal_dimension_vs_n method
experiment.plot_minimal_dimension_vs_n(k)

# 7. Call the plot_violations_vs_d method
experiment.plot_violations_vs_d(k)
\end{lstlisting}